% Navigation zwischen der Ganzzahlaxiomatik und ihren beiden Ergänzungen.
\begin{luacode*}
  local job = tex.jobname or ""
  local directory, main_directory = "registry/integers/", "registry/"
  if job:match("^_B[0-9][0-9]$") then
    directory, main_directory = "integers/", ""
  elseif job:match("^_B[0-9][0-9]%-") then
    directory, main_directory = "../integers/", "../"
  end
  if job:match("^_B17%-integers%-") then directory = "" end
  tex.sprint("\\gdef\\IntegersEditionDirectory{" .. directory .. "}")
  tex.sprint("\\gdef\\IntegersMainDirectory{" .. main_directory .. "}")
\end{luacode*}
% Private Nummern der getrennten Konstruktion.
\expandafter\def\csname IntegersPrivate@IntAddAssociativeOnClasses\endcsname{Th.~17E.2.4.11(H1)}
\expandafter\def\csname IntegersPrivate@IntAddClosureForRepresentatives\endcsname{Th.~17E.2.4.7(H1)}
\expandafter\def\csname IntegersPrivate@IntAddClosureRepresentativeElimination\endcsname{Th.~17E.2.4.7(H2)}
\expandafter\def\csname IntegersPrivate@IntAddCommutativeOnRepresentatives\endcsname{Th.~17E.2.4.12(H1)}
\expandafter\def\csname IntegersPrivate@IntAddCommutativeRepresentativeElimination\endcsname{Th.~17E.2.4.12(H2)}
\expandafter\def\csname IntegersPrivate@IntAddInverseOnClass\endcsname{Th.~17E.2.4.14(H1)}
\expandafter\def\csname IntegersPrivate@IntAddRightZeroForRepresentative\endcsname{Th.~17E.2.4.13(H1)}
\expandafter\def\csname IntegersPrivate@IntAddZeroRepresentativeElimination\endcsname{Th.~17E.2.4.13(H2)}
\expandafter\def\csname IntegersPrivate@IntMulAssociativeNegativeCoordinate\endcsname{Th.~17E.2.5.6(H8)}
\expandafter\def\csname IntegersPrivate@IntMulAssociativeNegativeExpandedComparison\endcsname{Th.~17E.2.5.6(H7)}
\expandafter\def\csname IntegersPrivate@IntMulAssociativeNegativeLeftExpansion\endcsname{Th.~17E.2.5.6(H5)}
\expandafter\def\csname IntegersPrivate@IntMulAssociativeNegativeRightExpansion\endcsname{Th.~17E.2.5.6(H6)}
\expandafter\def\csname IntegersPrivate@IntMulAssociativeOnClasses\endcsname{Th.~17E.2.5.6(H9)}
\expandafter\def\csname IntegersPrivate@IntMulAssociativePositiveCoordinate\endcsname{Th.~17E.2.5.6(H4)}
\expandafter\def\csname IntegersPrivate@IntMulAssociativePositiveExpandedComparison\endcsname{Th.~17E.2.5.6(H3)}
\expandafter\def\csname IntegersPrivate@IntMulAssociativePositiveLeftExpansion\endcsname{Th.~17E.2.5.6(H1)}
\expandafter\def\csname IntegersPrivate@IntMulAssociativePositiveRightExpansion\endcsname{Th.~17E.2.5.6(H2)}
\expandafter\def\csname IntegersPrivate@IntMulClosureForRepresentatives\endcsname{Th.~17E.2.5.3(H1)}
\expandafter\def\csname IntegersPrivate@IntMulClosureRepresentativeElimination\endcsname{Th.~17E.2.5.3(H2)}
\expandafter\def\csname IntegersPrivate@IntMulCommutativeClassCoordinates\endcsname{Th.~17E.2.5.4(H2)}
\expandafter\def\csname IntegersPrivate@IntMulCommutativeRepresentativeElimination\endcsname{Th.~17E.2.5.4(H3)}
\expandafter\def\csname IntegersPrivate@IntMulCommutativeRepresentativeProducts\endcsname{Th.~17E.2.5.4(H1)}
\expandafter\def\csname IntegersPrivate@IntMulLeftDistributiveNegativeCoordinate\endcsname{Th.~17E.2.5.7(H2)}
\expandafter\def\csname IntegersPrivate@IntMulLeftDistributiveOnClasses\endcsname{Th.~17E.2.5.7(H3)}
\expandafter\def\csname IntegersPrivate@IntMulLeftDistributivePositiveCoordinate\endcsname{Th.~17E.2.5.7(H1)}
\expandafter\def\csname IntegersPrivate@IntMulRightIdentityRepresentative\endcsname{Th.~17E.2.5.5(H1)}
\expandafter\def\csname IntegersPrivate@IntMulTwoSidedIdentityConclusion\endcsname{Th.~17E.2.5.5(H2)}
\expandafter\def\csname IntegersPrivate@IntPairRelDef\endcsname{Def.~17E.2.1.2}
\expandafter\def\csname IntegersPrivate@IntRelationReflexivePairs\endcsname{Th.~17E.2.1.1}
\expandafter\def\csname IntegersPrivate@IntRelationSymmetricPairs\endcsname{Th.~17E.2.1.2}
\expandafter\def\csname IntegersPrivate@IntRelationTransitiveCancellation\endcsname{Th.~17E.2.1.3(H2)}
\expandafter\def\csname IntegersPrivate@IntRelationTransitivePairs\endcsname{Th.~17E.2.1.3}
\expandafter\def\csname IntegersPrivate@IntRelationTransitiveRearrangement\endcsname{Th.~17E.2.1.3(H1)}
\expandafter\def\csname IntegersPrivate@IntegerAddOneClosure\endcsname{Th.~17E.2.4.9}
\expandafter\def\csname IntegersPrivate@IntegerAdditionCompatible\endcsname{Th.~17E.2.4.5}
\expandafter\def\csname IntegersPrivate@IntegerAdditionInputCrossSums\endcsname{Th.~17E.2.4.5(H1)}
\expandafter\def\csname IntegersPrivate@IntegerAdditionOutputCrossSum\endcsname{Th.~17E.2.4.5(H2)}
\expandafter\def\csname IntegersPrivate@IntegerAdditionQuotientDescent\endcsname{Th.~17E.2.4.6}
\expandafter\def\csname IntegersPrivate@IntegerClassDef\endcsname{Def.~17E.2.2.2}
\expandafter\def\csname IntegersPrivate@IntegerClassEqualityViaRelation\endcsname{Th.~17E.2.2.2(H1)}
\expandafter\def\csname IntegersPrivate@IntegerLeqStrictTransitiveFromAxioms\endcsname{Th.~17E.4.2.19}
\expandafter\def\csname IntegersPrivate@IntegerMinusOneClosure\endcsname{Th.~17E.2.4.10}
\expandafter\def\csname IntegersPrivate@IntegerMultiplicationCompatible\endcsname{Th.~17E.2.5.1}
\expandafter\def\csname IntegersPrivate@IntegerMultiplicationCompatibleLeft\endcsname{Th.~17E.2.5.1(H1)}
\expandafter\def\csname IntegersPrivate@IntegerMultiplicationCompatibleRight\endcsname{Th.~17E.2.5.1(H2)}
\expandafter\def\csname IntegersPrivate@IntegerMultiplicationQuotientDescent\endcsname{Th.~17E.2.5.2}
\expandafter\def\csname IntegersPrivate@IntegerNegationCompatible\endcsname{Th.~17E.2.4.1}
\expandafter\def\csname IntegersPrivate@IntegerNegationCompatibleCrossSum\endcsname{Th.~17E.2.4.1(H1)}
\expandafter\def\csname IntegersPrivate@IntegerNegationQuotientDescent\endcsname{Th.~17E.2.4.2}
\expandafter\def\csname IntegersPrivate@IntegerNegativeZero\endcsname{Th.~17E.2.7.2}
\expandafter\def\csname IntegersPrivate@IntegerNormalFormNonnegative\endcsname{Th.~17E.3.2.3(H1)}
\expandafter\def\csname IntegersPrivate@IntegerNormalFormNonpositive\endcsname{Th.~17E.3.2.3(H2)}
\expandafter\def\csname IntegersPrivate@IntegerNormalFormSuccessorBound\endcsname{Th.~17E.4.4.2}
\expandafter\def\csname IntegersPrivate@IntegerNumeralConvention\endcsname{Def.~17E.2.3.3}
\expandafter\def\csname IntegersPrivate@IntegerOnePositiveModel\endcsname{Th.~17E.4.2.7}
\expandafter\def\csname IntegersPrivate@IntegerOrderAntisymmetric\endcsname{Th.~17E.4.2.3(H2)}
\expandafter\def\csname IntegersPrivate@IntegerOrderComparable\endcsname{Th.~17E.4.2.3(H4)}
\expandafter\def\csname IntegersPrivate@IntegerOrderCompatible\endcsname{Th.~17E.4.2.1}
\expandafter\def\csname IntegersPrivate@IntegerOrderCompatibleBackward\endcsname{Th.~17E.4.2.1(H2)}
\expandafter\def\csname IntegersPrivate@IntegerOrderCompatibleForward\endcsname{Th.~17E.4.2.1(H1)}
\expandafter\def\csname IntegersPrivate@IntegerOrderDef\endcsname{Def.~17E.4.2.1}
\expandafter\def\csname IntegersPrivate@IntegerOrderNotation\endcsname{Def.~17E.4.2.2}
\expandafter\def\csname IntegersPrivate@IntegerOrderReflexive\endcsname{Th.~17E.4.2.3(H1)}
\expandafter\def\csname IntegersPrivate@IntegerOrderTransitive\endcsname{Th.~17E.4.2.3(H3)}
\expandafter\def\csname IntegersPrivate@IntegerOrderTranslationModel\endcsname{Th.~17E.4.2.4}
\expandafter\def\csname IntegersPrivate@IntegerOrderTranslationOnClasses\endcsname{Th.~17E.4.2.4(H1)}
\expandafter\def\csname IntegersPrivate@IntegerPairCarrier\endcsname{Def.~17E.2.1.1}
\expandafter\def\csname IntegersPrivate@IntegerPairClassTyping\endcsname{Th.~17E.2.2.4}
\expandafter\def\csname IntegersPrivate@IntegerPairEquivalence\endcsname{Th.~17E.2.1.4}
\expandafter\def\csname IntegersPrivate@IntegerPairEquivalenceReflexive\endcsname{Th.~17E.2.1.4(H1)}
\expandafter\def\csname IntegersPrivate@IntegerPairEquivalenceSymmetric\endcsname{Th.~17E.2.1.4(H2)}
\expandafter\def\csname IntegersPrivate@IntegerPairEquivalenceTransitive\endcsname{Th.~17E.2.1.4(H3)}
\expandafter\def\csname IntegersPrivate@IntegerPairRepresentationSurjective\endcsname{Th.~17E.2.2.5}
\expandafter\def\csname IntegersPrivate@IntegerPeanoModelCarrierMaps\endcsname{Th.~17E.5.2.1(H1)}
\expandafter\def\csname IntegersPrivate@IntegerPeanoModelInduction\endcsname{Th.~17E.5.2.1(H4)}
\expandafter\def\csname IntegersPrivate@IntegerPeanoModelInverses\endcsname{Th.~17E.5.2.1(H2)}
\expandafter\def\csname IntegersPrivate@IntegerPeanoModelOrder\endcsname{Th.~17E.5.2.1(H3)}
\expandafter\def\csname IntegersPrivate@IntegerPositiveEmbeddingRepresentativeModel\endcsname{Th.~17E.4.2.8}
\expandafter\def\csname IntegersPrivate@IntegerPositiveEmbeddingRepresentativeNegativeCase\endcsname{Th.~17E.4.2.8(H2)}
\expandafter\def\csname IntegersPrivate@IntegerPositiveEmbeddingRepresentativePositiveCase\endcsname{Th.~17E.4.2.8(H1)}
\expandafter\def\csname IntegersPrivate@IntegerPositiveFactorOrder\endcsname{Th.~17E.4.3.5}
\expandafter\def\csname IntegersPrivate@IntegerPositiveFactorOrderModel\endcsname{Th.~17E.4.2.16}
\expandafter\def\csname IntegersPrivate@IntegerPositiveFactorOrderOnRepresentativesModel\endcsname{Th.~17E.4.2.16(H1)}
\expandafter\def\csname IntegersPrivate@IntegerPositiveProductModel\endcsname{Th.~17E.4.2.12}
\expandafter\def\csname IntegersPrivate@IntegerPositiveProductNonnegativeFromAxioms\endcsname{Th.~17E.4.2.13}
\expandafter\def\csname IntegersPrivate@IntegerPredAfterSucc\endcsname{Th.~17E.3.1.7}
\expandafter\def\csname IntegersPrivate@IntegerPredAfterSuccExpansion\endcsname{Th.~17E.3.1.7(H1)}
\expandafter\def\csname IntegersPrivate@IntegerPredecessorClassValue\endcsname{Th.~17E.3.1.6}
\expandafter\def\csname IntegersPrivate@IntegerPredecessorDef\endcsname{Def.~17E.3.1.2}
\expandafter\def\csname IntegersPrivate@IntegerPredecessorFunction\endcsname{Th.~17E.3.1.2}
\expandafter\def\csname IntegersPrivate@IntegerPredecessorOnNegativeEmbedding\endcsname{Th.~17E.3.2.2}
\expandafter\def\csname IntegersPrivate@IntegerPredecessorValue\endcsname{Th.~17E.3.1.4}
\expandafter\def\csname IntegersPrivate@IntegerQuotientRepresentative\endcsname{Th.~17E.2.2.3(H1)}
\expandafter\def\csname IntegersPrivate@IntegerStrictLeqTransitiveFromAxioms\endcsname{Th.~17E.4.2.20}
\expandafter\def\csname IntegersPrivate@IntegerSubtractionClosure\endcsname{Th.~17E.2.4.8}
\expandafter\def\csname IntegersPrivate@IntegerSubtractionDef\endcsname{Def.~17E.2.4.3}
\expandafter\def\csname IntegersPrivate@IntegerSuccAfterPred\endcsname{Th.~17E.3.1.8}
\expandafter\def\csname IntegersPrivate@IntegerSuccAfterPredConclusion\endcsname{Th.~17E.3.1.8(H2)}
\expandafter\def\csname IntegersPrivate@IntegerSuccAfterPredExpansion\endcsname{Th.~17E.3.1.8(H1)}
\expandafter\def\csname IntegersPrivate@IntegerSuccessorBijective\endcsname{Th.~17E.3.1.9}
\expandafter\def\csname IntegersPrivate@IntegerSuccessorClassValue\endcsname{Th.~17E.3.1.5}
\expandafter\def\csname IntegersPrivate@IntegerSuccessorDef\endcsname{Def.~17E.3.1.1}
\expandafter\def\csname IntegersPrivate@IntegerSuccessorDistinct\endcsname{Th.~17E.4.4.1(H2)}
\expandafter\def\csname IntegersPrivate@IntegerSuccessorFunction\endcsname{Th.~17E.3.1.1}
\expandafter\def\csname IntegersPrivate@IntegerSuccessorInjective\endcsname{Th.~17E.3.1.9(H1)}
\expandafter\def\csname IntegersPrivate@IntegerSuccessorLeq\endcsname{Th.~17E.4.4.1(H1)}
\expandafter\def\csname IntegersPrivate@IntegerSuccessorPositive\endcsname{Th.~17E.4.4.1}
\expandafter\def\csname IntegersPrivate@IntegerSuccessorSurjective\endcsname{Th.~17E.3.1.9(H2)}
\expandafter\def\csname IntegersPrivate@IntegerSuccessorValue\endcsname{Th.~17E.3.1.3}
\expandafter\def\csname IntegersPrivate@IntegerTwoSidedInductionModel\endcsname{Th.~17E.5.1.1}
\expandafter\def\csname IntegersPrivate@IntegerTwoSidedInductionNegativeRay\endcsname{Th.~17E.5.1.1(H2)}
\expandafter\def\csname IntegersPrivate@IntegerTwoSidedInductionPositiveRay\endcsname{Th.~17E.5.1.1(H1)}
\expandafter\def\csname IntegersPrivate@IntegerZeroNeOneModel\endcsname{Th.~17E.4.2.6}
\expandafter\def\csname IntegersPrivate@IntegersAsQuotient\endcsname{Def.~17E.2.2.1}
\expandafter\def\csname IntegersPrivate@NaturalCopyInIntegers\endcsname{Def.~17E.2.3.2}
\expandafter\def\csname IntegersPrivate@NaturalCopyIntegerBijection\endcsname{Th.~17E.2.3.5}
\expandafter\def\csname IntegersPrivate@NaturalIntegerEmbedding\endcsname{Def.~17E.2.3.1}
\expandafter\def\csname IntegersPrivate@NaturalIntegerEmbeddingMultiplicative\endcsname{Th.~17E.2.6.2}
\expandafter\def\csname IntegersPrivate@NaturalIntegerEmbeddingMultiplicativeConclusion\endcsname{Th.~17E.2.6.2(H2)}
\expandafter\def\csname IntegersPrivate@NaturalIntegerEmbeddingOrderPreserving\endcsname{Th.~17E.4.2.9}
\expandafter\def\csname IntegersPrivate@NaturalIntegerEmbeddingPositiveModel\endcsname{Th.~17E.4.2.10}
\expandafter\def\csname IntegersPrivate@NaturalIntegerEmbeddingProductClass\endcsname{Th.~17E.2.6.2(H1)}
\expandafter\def\csname IntegersPrivate@NaturalIntegerEmbeddingReflectsEquality\endcsname{Th.~17E.2.3.3(H1)}
\expandafter\def\csname IntegersPrivate@NaturalNonzeroProductForIntegerModel\endcsname{Th.~17E.4.2.11}
\expandafter\def\csname IntegersPrivate@NaturalOrderAdditionFirstTranslation\endcsname{Th.~17E.4.1.2(H1)}
\expandafter\def\csname IntegersPrivate@NaturalOrderAdditionForIntegers\endcsname{Th.~17E.4.1.2}
\expandafter\def\csname IntegersPrivate@NaturalOrderAdditionSecondTranslation\endcsname{Th.~17E.4.1.2(H2)}
\expandafter\def\csname IntegersPrivate@NaturalOrderAdditionTransitiveConclusion\endcsname{Th.~17E.4.1.2(H3)}
\expandafter\def\csname IntegersPrivate@NaturalOrderTranslationBackward\endcsname{Th.~17E.4.1.1(H2)}
\expandafter\def\csname IntegersPrivate@NaturalOrderTranslationForIntegers\endcsname{Th.~17E.4.1.1}
\expandafter\def\csname IntegersPrivate@NaturalOrderTranslationForward\endcsname{Th.~17E.4.1.1(H1)}
\expandafter\def\csname IntegersPrivate@NaturalPositiveFactorOrderBackwardForIntegerModel\endcsname{Th.~17E.4.2.15(H2)}
\expandafter\def\csname IntegersPrivate@NaturalPositiveFactorOrderForIntegerModel\endcsname{Th.~17E.4.2.15}
\expandafter\def\csname IntegersPrivate@NaturalPositiveFactorOrderForwardForIntegerModel\endcsname{Th.~17E.4.2.15(H1)}
\expandafter\def\csname IntegersPrivate@NaturalPositiveFactorStrictOrderForIntegerModel\endcsname{Th.~17E.4.2.14}
\expandafter\def\csname IntegersPrivate@PositiveIntegerEmbeddingFactorLowerBound\endcsname{Th.~17E.4.3.7}
\expandafter\def\csname IntegersPrivate@PositiveIntegerOneLowerBound\endcsname{Th.~17E.4.3.3}

\newcommand{\IntegersReadingLink}{%
  \href{\IntegersEditionDirectory_B17-integers-reading.pdf\#integers.reading}{Lesefassung zu den ganzen Zahlen}%
}
\newcommand{\IntegersProofsLink}{%
  \href{\IntegersEditionDirectory_B17-integers-proofs.pdf\#integers.proofs}{Beweistabellen zu den ganzen Zahlen}%
}
\newcommand{\IntegersMainLink}{%
  \href{\IntegersMainDirectory_B17.pdf\#integers.main}{Modellnachweis in Band~17}%
}
\newcommand{\IntegersProofFor}[1]{%
  \href{\IntegersEditionDirectory_B17-integers-proofs.pdf\#integers.proof.#1}{vollständige Ableitung im Ergänzungsband}%
}
\NewDocumentCommand{\IntegersResultRef}{m o}{%
  \ifcsname IntegersPrivate@#1\endcsname
    \href{\IntegersEditionDirectory_B17-integers-proofs.pdf\#integers.statement.#1}%
      {\csname IntegersPrivate@#1\endcsname}%
  \else
    \IfNoValueTF{#2}{\FormulaRefAuto{#1}}{\FormulaRefAuto{#1}[#2]}%
  \fi
}
